\documentclass[UTF8,fontset=fandol,11pt]{ctexart}
\usepackage[a4paper,margin=1.75cm]{geometry}
\usepackage{amsmath,amssymb}
\usepackage{enumitem}
\usepackage{longtable}
\usepackage{booktabs}
\usepackage{tcolorbox}
\tcbuselibrary{skins,breakable}
\usepackage{xcolor}
\usepackage{hyperref}
\usepackage{graphicx}
\usepackage{tabularx}
\usepackage{fancyhdr}
\usepackage{titlesec}
\usepackage{array}
\hypersetup{colorlinks=true,linkcolor=blue,urlcolor=blue}
\pagestyle{fancy}
\fancyhf{}
\lhead{\small 工程管理复习讲义}
\rhead{\small \thepage}
\setlength{\parindent}{0pt}
\setlist[enumerate]{leftmargin=2em,itemsep=0.3em}
\setlist[itemize]{leftmargin=1.8em,itemsep=0.2em}
\definecolor{mainblue}{HTML}{2454D6}
\definecolor{deepblue}{HTML}{13204B}
\definecolor{softblue}{HTML}{EEF5FF}
\definecolor{mainpurple}{HTML}{7B2CBF}
\definecolor{softpurple}{HTML}{F5ECFF}
\definecolor{mainorange}{HTML}{F77F00}
\definecolor{softorange}{HTML}{FFF3E0}
\definecolor{maingreen}{HTML}{138A36}
\definecolor{softgreen}{HTML}{EAF8ED}
\definecolor{answerred}{HTML}{D62828}
\newtcolorbox{mainbox}[1][]{enhanced,colback=softblue,colframe=mainblue,boxrule=0.9pt,arc=2mm,left=2.5mm,right=2.5mm,top=1.5mm,bottom=1.5mm,#1}
\newtcolorbox{quizbox}[1][]{enhanced,breakable,colback=white,colframe=mainpurple,coltitle=white,colbacktitle=mainpurple,fonttitle=\bfseries,boxrule=0.8pt,arc=2mm,left=2.8mm,right=2.8mm,top=1.6mm,bottom=1.6mm,#1}
\newtcolorbox{answerbox}[1][]{enhanced,breakable,colback=softgreen,colframe=maingreen,coltitle=white,colbacktitle=maingreen,fonttitle=\bfseries,boxrule=0.7pt,arc=2mm,left=2.5mm,right=2.5mm,top=1mm,bottom=1mm,#1}
\newtcolorbox{reviewbox}[1][]{enhanced,breakable,colback=softorange,colframe=mainorange,coltitle=white,colbacktitle=mainorange,fonttitle=\bfseries,boxrule=0.7pt,arc=2mm,left=2.5mm,right=2.5mm,top=1mm,bottom=1mm,#1}
\newcommand{\chaptertitle}[1]{\vspace{0.5em}\begin{tcolorbox}[enhanced,colback=deepblue,colframe=deepblue,arc=2mm,boxrule=0pt]\begin{center}\Large\bfseries\color{white}#1\end{center}\end{tcolorbox}\vspace{0.5em}}
\newcommand{\sectitle}[1]{\vspace{0.5em}\noindent{\Large\bfseries\color{mainblue}#1}\par\vspace{0.25em}{\color{mainorange}\hrule height 1.2pt}\vspace{0.45em}}
\newcommand{\answerline}{\textcolor{answerred}{\rule{2.4cm}{0.7pt}}}

\begin{document}

\tableofcontents

\newpage

\chaptertitle{PPT 重点提要}

\section*{1. 危险性较大的分部分项工程安全管理规定}

\addcontentsline{toc}{section}{1. 危险性较大的分部分项工程安全管理规定}

\begin{mainbox}

\textbf{提要：} 以下内容按课件自动提炼，适合快速背诵。

\end{mainbox}

\begin{enumerate}[label=\arabic*.]

\item 二、危险性较大的分部分项工程安全管理规定

\item 前期保障

\item （

\item ）建设单位应当组织勘察、设计等单位在施工招标文件中列出危大工程清单，要求施工单位在投标时补充完善危大工程清单并明确相应的安全管理措施。

\item ）建设单位应当按照施工合同约定及时支付危大工程施工技术措施费以及相应的安全防护文明施工措施费，保障危大工程施工安全。

\item ）建设单位在申请办理安全监督手续时，应当提交危大工程清单及其安全管理措施等资料。

\item 专项施工方案

\item ）施工单位应当在危大工程施工前组织工程技术人员编制专项施工方案。

\item ）实行施工总承包的，专项施工方案应当由施工总承包单位组织编制。其中，起重机械安装拆卸工程、深基坑工程、附着式升降脚手架等专业工程实行分包的，其专项施工方案可以由相关专业分包单位组织编制。

\item ）专项施工方案应当由施工单位技术负责人审核签字、加盖单位公章，并由总监理工程师审查签字、加盖执业印章后方可实施。

\item ）危大工程实行分包并由分包单位编制专项施工方案的，专项施工方案应当由总承包单位技术负责人及分包单位技术负责人共同审核签字并加盖单位公章。

\item ）专家应当从地方人民政府住房城乡建设主管部门建立的专家库中选取，符合专业要求且人数不得少于

\item 名。与本工程有利害关系的人员不得以专家身份参加专家论证会。

\item ）专家论证会后，应当形成论证报告，对专项施工方案提出通过、修改后通过或者不通过的一致意见。专家对论证报告负责并签字确认。

\item ）专项施工方案经论证需修改后通过的，施工单位应当根据论证报告修改完善后，重新履行审核（查）、签字、盖章的程序。

\item ）专项施工方案经论证不通过的，施工单位修改后应当按照本规定的要求重新组织专家论证。

\item ）对于超过一定规模的危大工程，施工单位应当组织召开专家论证会对专项施工方案进行论证。实行施工总承包的，由施工总承包单位组织召开专家论证会。专家论证前专项施工方案应当通过施工单位审核和总监理工程师审查。

\item 现场安全管理

\item ）施工单位应当在施工现场显著位置公告危大工程名称、施工时间和具体责任人员，并在危险区域设置安全警示标志。

\item ）专项施工方案实施前，编制人员或者项目技术负责人应当向施工现场管理人员进行方案交底。

\item ）施工现场管理人员应当向作业人员进行 安全技术交底，并由双方和项目专职安全生产管理人员共同签字确认。

\item ）施工单位应当严格按照专项施工方案组织施工，不得擅自修改专项施工方案。

\item ）因规划调整、设计变更等原因确需调整的，修改后的专项施工方案应当按照本规定重新审核和论证。涉及资金或者工期调整的，建设单位应当按照约定予以调整。

\item ）施工单位应当对危大工程施工作业人员进行登记，项目负责人应当在施工现场履职。

\end{enumerate}

\newpage

\section*{2. 双代号时标网络计划}

\addcontentsline{toc}{section}{2. 双代号时标网络计划}

\begin{mainbox}

\textbf{提要：} 以下内容按课件自动提炼，适合快速背诵。

\end{mainbox}

\begin{enumerate}[label=\arabic*.]

\item 一、双代号时标网络计划

\item （一）双代号时标网络计划的特点及其规定

\item 、双代号时标网络计划的概念

\item 以时间坐标为尺度表示工作时间的网络计划。

\item 、双代号时标网络计划的特点

\item 、双代号时标网络计划的规定

\item 一般规定：

\item （

\item ）实箭线表示一项工作，开始节点与结束节点中心的水平距离表示该工作的持续时间。

\item ）虚箭线表示一项虚工作。

\item 因为其持续时间为零，所以虚箭线应竖向绘制。

\item ）波形线表示一项工作的自由时差。

\item 注意：

\item ）波形线应绘制在水平实线之后。

\item ）当一项虚工作存在

\item FF

\item ij

\item 时，按下述方法绘制：

\item （二）双代号时标网络计划的绘制方法

\item 、绘制要求：

\item 绘制要求

\item ）时标网络计划需绘制在带有时间坐标的表格上；

\item ）节点中心必须对准时间坐标的刻度线；

\item ）实箭线表示实工作、虚箭线表示虚工作、水平破浪线表示自由时差；

\end{enumerate}

\newpage

\section*{3. 双代号网络计划1}

\addcontentsline{toc}{section}{3. 双代号网络计划1}

\begin{mainbox}

\textbf{提要：} 以下内容按课件自动提炼，适合快速背诵。

\end{mainbox}

\begin{enumerate}[label=\arabic*.]

\item 13.1

\item 网络计划原理

\item （一）概述

\item 世纪

\item 年代

\item 美国杜邦公司

\item 由我国著名数学家华罗庚引入并创新了统筹法

\item 统筹法

\item 网络计划

\item 广义的

\item 狭义的

\item 其应用不仅仅局限于工程项目，规模庞大、管理复杂的项目均可采用，美国北极星导弹潜艇的研制、阿波罗登月计划；

\item 以工程项目为对象编制的，全称为工程网络计划，特别是对于庞大的工程项目，如果网络计划运用得当，不仅能够方便控制工期，还能挖潜增效、节约资源、降低成本；

\item （二）网络图

\item 由

\item 箭线

\item 和

\item 节点

\item 组成的，用来表示工作流程的

\item 有向

\item 、

\item 有序

\item 的网状图形。

\item 其中有向有序是指遵循网络图的

\end{enumerate}

\newpage

\section*{4. 双代号网络计划2}

\addcontentsline{toc}{section}{4. 双代号网络计划2}

\begin{mainbox}

\textbf{提要：} 以下内容按课件自动提炼，适合快速背诵。

\end{mainbox}

\begin{enumerate}[label=\arabic*.]

\item （九）双代号网络计划时间参数的计算

\item 、时间参数计算的目的

\item 、时间参数

\item ——

\item “

\item 整体

\item 参数”

\item 工作

\item 、时间参数的计算方法

\item “工作计算法”（六时标注法）

\item （

\item ）最早（开始

\item /

\item 完成）时间的计算

\item ）最迟（完成

\item 开始）时间的计算

\item ）计算自由时差

\item ）计算总时差

\item ）确定关键工作和关键线路

\item “节点计算法”（二时标注法）

\item 节点最早时间计算

\item 节点最迟时间计算

\item 当计划工期等于计算工期时，

\item 关键节点的最早时间与最迟时间必然相等

\end{enumerate}

\newpage

\section*{5. 土方回填}

\addcontentsline{toc}{section}{5. 土方回填}

\begin{mainbox}

\textbf{提要：} 以下内容按课件自动提炼，适合快速背诵。

\end{mainbox}

\begin{enumerate}[label=\arabic*.]

\item 四、土方回填

\item （一）土料要求

\item 土料要求

\item （

\item ）填料应符合设计要求，保证填方的

\item 强度和稳定性

\item 。

\item ）

\item 不能选用

\item ：淤泥、淤泥质土、膨胀土、有机质>

\item 5\%

\item 的土、含水量不符合压实要求的黏性土。

\item ）填方土应尽量采用

\item 同类土

\item ）土料含水量一般以

\item 手握成团、落地开花

\item 为适宜。

\item （二）基底处理

\item 基底处理

\item 清除

\item 基底上的垃圾、草皮、树根、杂物，排除坑穴中积水、淤泥和种植土，将基底充分夯实和碾压密实。

\item ）应

\item 采取措施防止

\item 地表滞水流入填方区，浸泡地基，造成基土下陷。

\end{enumerate}

\newpage

\section*{6. 土方工程概述}

\addcontentsline{toc}{section}{6. 土方工程概述}

\begin{mainbox}

\textbf{提要：} 以下内容按课件自动提炼，适合快速背诵。

\end{mainbox}

\begin{enumerate}[label=\arabic*.]

\item 一、土方工程概述

\item （一）土方工程的内容

\item 包括一切土的挖掘、填筑、运输等过程以及排水降水、土壁支撑等准备工作和辅助工程。

\item （二）土的工程分类

\item 《

\item 水利水电工程施工组织设计规范

\item 》

\item （

\item SL 303-2017

\item ）：

\item 公路工程预算定额

\item JTG-T 3832-2018

\item （三）土的工程性质

\item 工程性质

\item 内容

\item 内摩擦角

\item 土体中颗粒间相互移动和胶合作用形成的摩擦特性。其数值为强度包线与水平线的夹角——

\item 土的抗剪强度指标

\item 。

\item 土抗剪强度

\item 土体抵抗剪切破坏的极限强度，包括

\item 内摩擦力和内聚力

\item 黏聚力

\item 同种物质内部相邻各部分之间的

\end{enumerate}

\newpage

\section*{7. 土方工程的机械化施工}

\addcontentsline{toc}{section}{7. 土方工程的机械化施工}

\begin{mainbox}

\textbf{提要：} 以下内容按课件自动提炼，适合快速背诵。

\end{mainbox}

\begin{enumerate}[label=\arabic*.]

\item 二、土方工程的机械化施工

\item （一）主要土方机械的特点与施工方法

\item 、推土机

\item 推土机可推挖一至三类土，其推运距离宜在

\item 100m

\item 以内，

\item 30\textasciitilde{}60m

\item 时经济效果最好。

\item 土方推送

\item 水中清渣

\item 移挖作填

\item 平整场地

\item 、铲运机

\item 拖式铲运机运距以

\item 800m

\item 以内为宜，

\item 300m

\item 左右时效率最高。自行式铲运机的经济运距为

\item 800m\textasciitilde{}1500m

\item 。

\item 、单斗挖掘机

\item （

\item ）正铲挖掘机

\item ）反铲挖掘机

\end{enumerate}

\newpage

\section*{8. 土方开挖}

\addcontentsline{toc}{section}{8. 土方开挖}

\begin{mainbox}

\textbf{提要：} 以下内容按课件自动提炼，适合快速背诵。

\end{mainbox}

\begin{enumerate}[label=\arabic*.]

\item 三、土方开挖

\item 土方开挖

\item 方案确定

\item 应考虑

\item 土方量、土方运距、土方施工顺序、地质条件等因素，进行土方平衡和调配。

\item 原则

\item 开槽支撑，先撑后挖，分层开挖，严禁超挖

\item 。

\item 基本要求

\item 严禁

\item 在基坑（槽）及建（构）筑物周边

\item 影响范围内

\item 堆放土方。

\item 基坑边界周围地面应设排水沟，对坡顶、坡面、坡脚

\item 采取降排水措施

\item （一）浅基坑的开挖

\item 浅基坑开挖

\item （

\item ）相邻基坑开挖时，应遵循

\item 先深后浅或同时进行

\item 的施工程序；挖土应

\item 自上而下、水平分段分层

\item 进行；

\item ）基坑开挖应尽量

\end{enumerate}

\newpage

\section*{9. 地下水处理}

\addcontentsline{toc}{section}{9. 地下水处理}

\begin{mainbox}

\textbf{提要：} 以下内容按课件自动提炼，适合快速背诵。

\end{mainbox}

\begin{enumerate}[label=\arabic*.]

\item 四、地下水处理

\item （一）排除地面水

\item 目的

\item 为了保证土方及后续工程施工的顺利进行，场地内的地面水和雨水应及时排走，以保持土体干燥。

\item 处理办法

\item 一般采用“疏”，“截”，“挡”的办法：

\item “疏”—排水沟；

\item “截”—截水沟；

\item “挡”—修筑土堤。

\item （二）降低地下水

\item 在

\item 《

\item 建筑地基基础工程施工规范

\item 》

\item （

\item GB51004-2015

\item ）中，第

\item 7.1.4

\item 条，依据场地的水文地质条件、基础规模、开挖深度、各土层的渗透性能等，可选择集水明排、降水以及回灌等方法单独或组合使用。常用地下水控制方法及适用条件宜符合下表的规定。

\item 、集水明排

\item 、地下水控制技术方案选择

\item 人工降水

\item 方案编制依据

\item ）工程地质、水文地质、周边环境条件；

\end{enumerate}

\newpage

\section*{10. 地基处理技术与桩的类型}

\addcontentsline{toc}{section}{10. 地基处理技术与桩的类型}

\begin{mainbox}

\textbf{提要：} 以下内容按课件自动提炼，适合快速背诵。

\end{mainbox}

\begin{enumerate}[label=\arabic*.]

\item 一、常用地基处理技术

\item （一）换填地基

\item 换填地基

\item （

\item ）换填厚度由

\item 设计确定

\item ，一般宜为

\item 0.5

\item ～

\item 3m

\item 。

\item ）换填地基压实标准要求：换填材料为

\item 灰土、粉煤灰

\item 时，压实系数为

\item $\geq$

\item 0.95

\item ；

\item 其他

\item 材料时，压实系数为

\item 0.97

\item ）换填地基施工时，

\item 不得

\item 在柱基、墙角及承重窗间墙下

\item 接缝

\end{enumerate}

\newpage

\section*{11. 基坑支护施工技术}

\addcontentsline{toc}{section}{11. 基坑支护施工技术}

\begin{mainbox}

\textbf{提要：} 以下内容按课件自动提炼，适合快速背诵。

\end{mainbox}

\begin{enumerate}[label=\arabic*.]

\item 二、基坑支护施工技术

\item 施工前应具备的资料

\item （

\item ）岩土工程勘察报告、施工所需的设计文件；

\item ）施工影响范围内的建（构）筑物、地下管网和障碍物资料、施工组织设计、专项施工方案和施工监测方案。

\item 基坑工程专项施工方案的内容

\item ）支护结构、地下水控制、土方开挖和回填等施工技术参数；

\item ）基坑工程施工工艺流程；

\item ）基坑工程施工方法；

\item ）基坑工程施工安全技术措施；

\item ）应急预案；

\item ）工程监测要求等。

\item （一）浅基坑的支护

\item 浅基坑的支护

\item ）斜柱支撑——适于开挖较大型、深度不大的基坑或使用机械挖土时。

\item ）锚拉支撑——适于开挖较大型、深度较深的基坑或使用机械挖土，不能安设横撑时使用。

\item ）型钢桩横挡板支撑——适于地下水位较低、深度不很大的一般黏性土层或砂土层中使用。

\item ）短桩横隔板支撑——适于开挖宽度大的基坑，当部分地段下部 放坡不够时使用。

\item ）临时挡土墙支撑——适于开挖宽度大的基坑，当部分地段下部放坡不够时使用。

\item ）挡土灌注桩支护——适用于开挖较大、较浅（<

\item 5m

\item ）基坑，邻近有建筑物，不允许背面地基有下沉、位移时采用。

\item ）叠袋式挡墙支护——适用于一般粘性土、面积大、开挖深度应在

\item 以内的浅基坑支护。

\end{enumerate}

\newpage

\section*{12. 基坑监测}

\addcontentsline{toc}{section}{12. 基坑监测}

\begin{mainbox}

\textbf{提要：} 以下内容按课件自动提炼，适合快速背诵。

\end{mainbox}

\begin{enumerate}[label=\arabic*.]

\item 三、基坑监测

\item 应实施监测的基坑工程

\item （

\item ）基坑设计安全等级为一、二级的基坑。

\item ）开挖深度$\geq$

\item 5m

\item 的下列基坑：

\item ）土质基坑；

\item ）极软岩基坑、破碎的软岩基坑、极破碎的岩体基坑；

\item ）上部为土体，下部为极软岩、破碎的软岩、极破碎的岩体构成的土岩组合基坑。

\item ）开挖深度

\item <5m

\item 但现场地质情况和周围环境较复杂的基坑工程。

\item 任务委托

\item 基坑工程施工前，应由建设方委托具备相应资质的第三方对基坑工程实施现场监测。

\item 基坑工程监测的总体要求

\item ）基坑工程施工前，应编制基坑工程监测方案。

\item ）应根据基坑工程安全等级、周边环境条件、支护类型及施工场地等确定基坑工程监测项目、监测点布置、监测方法、监测频率和监测预警。

\item ）至少进行围护墙顶部水平位移、沉降以及周边建筑、道路等沉降监测，并应根据项目技术设计条件对围护墙或土体深层水平位移、支护结构内力、土压力、孔隙水压力等进行监测。

\item ）监测点应沿基坑围护墙顶部周边布设，周边中部、阳角处应布点。

\item ）当基坑监测达到变形预警值，或基坑出现流砂、管涌、隆起、陷落，或基坑支护结构及周边环境出现大的变形时，应立即进行预警。

\item ）基坑降水应对水位降深进行监测，地下水回灌施工应对回灌量和水质进行监测。

\item ）逆作法施工应全过程进行监测。

\item 应进行专项论证的监测方案

\end{enumerate}

\newpage

\section*{13. 基坑验槽方法}

\addcontentsline{toc}{section}{13. 基坑验槽方法}

\begin{mainbox}

\textbf{提要：} 以下内容按课件自动提炼，适合快速背诵。

\end{mainbox}

\begin{enumerate}[label=\arabic*.]

\item 五、基坑验槽方法

\item （一）验槽具备的资料和条件

\item 验槽具备的资料和条件

\item （

\item ）

\item 勘察、设计、建设、监理、施工

\item 等相关单位技术人员到场；

\item ）地基基础

\item 设计文件

\item ；

\item ）岩土工程

\item 勘察报告

\item ）轻型动力

\item 触探记录（可不进行时除外）

\item ）地基处理或深基础施工质量

\item 检测报告

\item ）基底应为无扰动的原状土，

\item 留置有保护层时其厚度不应超过

\item 100mm

\item 。

\item （二）天然地基验槽

\item 天然地基验槽内容

\item ）根据勘察、设计文件核对基坑的位置、平面尺寸、坑底标高；

\item ）根据勘察报告核对坑底、坑边岩土体及地下水情况；

\end{enumerate}

\newpage

\section*{14. 基本概念}

\addcontentsline{toc}{section}{14. 基本概念}

\begin{mainbox}

\textbf{提要：} 以下内容按课件自动提炼，适合快速背诵。

\end{mainbox}

\begin{enumerate}[label=\arabic*.]

\item 一、流水施工原理

\item （一）流水施工的定义

\item （二）组织施工的方式

\item 思考：

\item 钢筋混凝土独立基础施工过程分为：支模板、绑钢筋、浇混凝土等工序，如果是一个独立基础，如何组织施工？如果是

\item 个独立基础，又该如何组织施工呢？

\item 施工的基本组织方式

\item 施工特点

\item 工艺流程

\item 资源利用

\item 平面或空间布置

\item 、顺序施工（依次施工）

\item （

\item ）概念：

\item 第一个施工段

\item 的所有施工过程全部施工完毕后，再进行

\item 第二个施工段

\item 的施工。

\item ）思考：

\item ）顺序施工的基本特点：

\item 、平行施工

\item ）概念：所有施工过程的各个施工段

\item 同时开工

\item ，

\end{enumerate}

\newpage

\section*{15. 填充墙砌体工程}

\addcontentsline{toc}{section}{15. 填充墙砌体工程}

\begin{mainbox}

\textbf{提要：} 以下内容按课件自动提炼，适合快速背诵。

\end{mainbox}

\begin{enumerate}[label=\arabic*.]

\item 四、填充墙砌体工程

\item 填充墙砌体工程

\item （

\item ）填充墙砌体工程通常采用烧结空心砖、蒸压加气混凝土砌块、轻骨料混凝土小型空心砌块等。

\item ）砌筑填充墙时，轻骨料混凝土小型空心砌块和蒸压加气混凝土砌块的产品龄期不应＜

\item 28d

\item ，蒸压加气混凝土砌块的含水率宜＜

\item 30\%

\item 。

\item ）烧结空心砖、蒸压加气混凝土砌块、轻骨料混凝土小型空心砌块等在运输、装卸过程中，严禁抛掷和倾倒。进场后应按品种、规格堆放整齐，堆置高度不宜超过

\item 2m

\item 。蒸压加气混凝土砌块在运输及堆放中应防止雨淋。

\item ）吸水率较小的轻骨料混凝土小型空心砌块及采用薄灰砌筑法施工的蒸压加气混凝土砌块，砌筑前不应对其浇（喷）水湿润。

\item ）轻骨料混凝土小型空心砌块或加气混凝土砌块墙如无切实有效措施，不得使用于下列部位：

\item ）建筑物防潮层以下部位；

\item ）长期浸水或化学侵蚀环境；

\item ）长期处于有振动源环境的墙体；

\item ）砌块表面经常处于

\item ℃以上的高温环境。

\item ）在 厨房、卫生间、浴室等处采用轻骨料混凝土小型空心砌块、蒸压加气混凝土砌块砌筑墙体时，墙底部宜现浇混凝土坎台，其高度应为

\item 150mm

\item ）填充墙拉结筋处的下皮小砌块宜用 半盲孔小砌块或用混凝土灌实孔洞的小砌块。薄灰砌筑法施工的蒸压加气混凝土砌块砌体，拉结筋应放置在砌块上表面设置的沟槽内。

\item ）蒸压加气混凝土砌块、轻骨料混凝土小型空心砌块不应与其他块体混砌，不同强度等级的同类块体也不得混砌。

\item ）加气混凝土墙上不得留设脚手眼。每一楼层内的砌块墙应连续砌完，不留接槎。如必须留槎时，应留斜槎。

\end{enumerate}

\newpage

\section*{16. 大体积混凝土工程}

\addcontentsline{toc}{section}{16. 大体积混凝土工程}

\begin{mainbox}

\textbf{提要：} 以下内容按课件自动提炼，适合快速背诵。

\end{mainbox}

\begin{enumerate}[label=\arabic*.]

\item 四、大体积混凝土工程

\item （一）施工组织设计

\item 施工组织设计

\item （

\item ）大体积混凝土施工应编制施工组织设计或施工技术方案，并应有环境保护和安全施工的技术措施。

\item ）主要内容包括：

\item ）大体积混凝土浇筑体温度应力和收缩应力计算结果；

\item ）施工阶段主要抗裂构造措施和温控指标的确定；

\item ）原材料优选、配合比设计、制备与运输计划；

\item ）主要施工设备和现场总平面布置；

\item ）温控监测设备和测试布置图；

\item ）浇筑顺序和施工进度计划

\item ;

\item ）保温和保湿养护方法；

\item ）应急预案和应急保障措施；

\item ）特殊部位和特殊气候条件下的施工措施。

\item （二）施工规定

\item 施工规定

\item ）大体积混凝土的设计强度等级宜为

\item C25

\item ～

\item C50,

\item 并可采用混凝土

\item 60d

\end{enumerate}

\newpage

\section*{17. 屋面防水工程施工}

\addcontentsline{toc}{section}{17. 屋面防水工程施工}

\begin{mainbox}

\textbf{提要：} 以下内容按课件自动提炼，适合快速背诵。

\end{mainbox}

\begin{enumerate}[label=\arabic*.]

\item 三、

\item 屋面防水工程施工

\item （一）屋面防水等级和设防要求

\item 屋面防水工程应根据建筑物的类别、重要程度、使用功能要求确定防水等级，并应按相应等级进行防水设防；对防水有特殊要求的建筑屋面，应进行专项防水设计。屋面防水等级和设防要求应符合下表的规定。

\item 防水等级

\item 建筑类别

\item 设防要求

\item I

\item 级

\item 重要建筑和高层建筑

\item 两道防水设防

\item II级

\item 一般建筑

\item 一道防水设防

\item （二）防水材料选择的基本原则

\item 防水材料选择的基本原则

\item （

\item ）外露使用的防水层，应选用耐紫外线、耐老化、耐候性好的防水材料；

\item ）上人屋面，应选用耐霉变、拉伸强度高的防水材料；

\item ）长期处于潮湿环境的屋面，应选用耐腐蚀、耐霉变、耐穿刺、耐长期水浸等性能的防水材料；

\item ）薄壳、装配式结构、钢结构及大跨度建筑屋面，应选用耐候性好、适应变形能力强的防水材料；

\item ）倒置式屋面应选用适应变形能力强、接缝密封保证率高的防水材料；

\item ）坡屋面应选用与基层粘结力强、感温性小的防水材料；

\item ）屋面接缝密封防水，应选用与基材粘结力强和耐候性好、适应位移能力强的密封材料；

\end{enumerate}

\newpage

\section*{18. 建筑防水材料的特性与应用}

\addcontentsline{toc}{section}{18. 建筑防水材料的特性与应用}

\begin{mainbox}

\textbf{提要：} 以下内容按课件自动提炼，适合快速背诵。

\end{mainbox}

\begin{enumerate}[label=\arabic*.]

\item 一、建筑防水材料的特性与应用

\item 建筑防水

\item 构造防水

\item 材料防水

\item 刚性防水

\item 柔性防水

\item 防水砂浆

\item 防水混凝土

\item 防水卷材

\item 堵漏灌浆材料

\item 防水涂料

\item 密封材料

\item （一）防水卷材

\item 防水卷材的分类

\item 改性沥青防水卷材

\item （

\item ）按照改性材料的不同分为：弹性体改性沥青防水卷材、塑性体改性沥青防水卷材和其他改性沥青防水卷材。

\item ）主要有弹性体

\item (SBS)

\item 改性沥青防水卷材、塑性体

\item (APP)

\item 改性沥青防水卷材、沥青复合胎柔性防水卷材、自粘橡胶改性沥青防水卷材、改性沥青聚乙烯胎防水卷材以及道桥用改性沥青防水卷材等。

\item ）

\item SBS

\end{enumerate}

\newpage

\section*{19. 抹灰工程施工}

\addcontentsline{toc}{section}{19. 抹灰工程施工}

\begin{mainbox}

\textbf{提要：} 以下内容按课件自动提炼，适合快速背诵。

\end{mainbox}

\begin{enumerate}[label=\arabic*.]

\item 二、抹灰工程施工

\item （一）抹灰工程分类与材料的技术要求

\item 分类

\item （

\item ）一般抹灰：水泥砂浆、水泥混合砂浆、聚合物水泥砂浆、粉刷石膏等抹灰。

\item ）保温层薄抹灰：保温层外面聚合物砂浆薄抹灰。

\item ）装饰抹灰：水刷石、斩假石、干粘石、假面砖等抹灰。

\item ）清水砌体勾缝：清水砌体砂浆勾缝、原浆勾缝。

\item 水泥

\item ）抹灰用的水泥其强度等级应$\geq$

\item 32.5MPa

\item ，砂浆的拉伸粘结强度、聚合物砂浆的保水率复验应合格。

\item ）白水泥和彩色水泥主要用于装饰抹灰；

\item ）不同品种、不同强度等级的水泥不得混用。

\item 砂子

\item 砂子宜选用中砂，砂子使用前应过筛（不大于

\item 5mm

\item 的筛孔），不得含有杂质；特细砂不宜使用。

\item 石灰膏

\item 抹灰用的石灰膏的熟化期不应少

\item 15d

\item ；石灰膏应细腻洁白，不得含有未熟化颗粒，已冻结风化的石灰膏不得使用。

\item 彩色石粒

\item 彩色石粒是由天然大理石破碎而成，具有多种颜色，多用作水磨石、水刷石、斩假石的骨料。

\end{enumerate}

\newpage

\section*{20. 拓展：装饰工程}

\addcontentsline{toc}{section}{20. 拓展：装饰工程}

\begin{mainbox}

\textbf{提要：} 以下内容按课件自动提炼，适合快速背诵。

\end{mainbox}

\begin{enumerate}[label=\arabic*.]

\item 第

\item 章 装饰工程

\item 一、建筑装饰装修材料

\item （

\item 一）饰面石材

\item 天然花岗岩

\item 特性

\item ）构造致密、强度高、密度大、吸水率极低、质地坚硬、耐磨，属酸性硬石材。

\item ）

\item SiO

\item 的含量常为

\item 60\%

\item 以上，其耐酸、抗风化、耐久性好，使用年限长。

\item ）花岗石所含石英在高温下会发生晶变，体积膨胀而开裂，不耐火。

\item 分类、等级及技术要求

\item ）分类：

\item ）按形状：毛光板（

\item MG

\item ）、普型板（

\item PX

\item ）、圆弧板（

\item HM

\item ）、异型板（

\item YX

\end{enumerate}

\newpage

\section*{21. 拓展：防水工程}

\addcontentsline{toc}{section}{21. 拓展：防水工程}

\begin{mainbox}

\textbf{提要：} 以下内容按课件自动提炼，适合快速背诵。

\end{mainbox}

\begin{enumerate}[label=\arabic*.]

\item 第

\item 章 防水工程

\item 一

\item 、建筑防水材料的特性与应用

\item 建筑防水

\item 构造

\item 防水

\item 材料

\item 刚性

\item 柔性

\item 防水砂浆

\item 防水混凝土

\item 防水卷材

\item 堵漏灌浆材料

\item 防水涂料

\item 密封材料

\item （

\item 一）防水卷材

\item 防水卷材的分类

\item 改性沥青防水卷材

\item ）按照改性材料的不同分为：弹性体改性沥青防水卷材、塑性体改性沥青防水卷材和其他改性沥青防水卷材。

\item ）主要有弹性体

\item (SBS)

\item 改性沥青防水卷材、塑性体

\end{enumerate}

\newpage

\section*{22. 有节奏流水施工}

\addcontentsline{toc}{section}{22. 有节奏流水施工}

\begin{mainbox}

\textbf{提要：} 以下内容按课件自动提炼，适合快速背诵。

\end{mainbox}

\begin{enumerate}[label=\arabic*.]

\item （七）流水施工基本组织方式

\item 流水施工

\item 非（无）节奏流水施工

\item 有节奏

\item 等节奏（固定节拍）流水施工

\item 异节奏流水施工

\item 异步距异节奏流水施工（一般的成倍节拍流水施工）

\item 等步距异节奏流水施工（加快的成倍节拍流水施工）

\item 、等节奏流水施工（固定节拍流水施工）

\item （

\item ）概念：是指参与流水施工的施工过程流水节拍彼此相等的流水施工组织方式，即

\item 同一施工过程在不同的施工段上流水节拍相等

\item ，

\item 不同的施工过程在同一施工段上的流水节拍也相等

\item 的流水施工方式。

\item ）基本特征：

\item ）主要参数的确定

\item 例题：

\item 某工程由

\item A

\item 、

\item B

\item C

\item D

\end{enumerate}

\newpage

\section*{23. 桩基检测技术}

\addcontentsline{toc}{section}{23. 桩基检测技术}

\begin{mainbox}

\textbf{提要：} 以下内容按课件自动提炼，适合快速背诵。

\end{mainbox}

\begin{enumerate}[label=\arabic*.]

\item 五、桩基检测技术

\item （一）检测阶段、作用及指标

\item 检测阶段

\item 检测作用

\item 检测指标

\item 施工前

\item 为设计提供依据的试验桩

\item 单桩极限承载力

\item 施工后

\item 为验收提供依据的工程桩

\item 单桩承载力

\item 桩身完整性检测

\item （二）桩基检测的方法和目的

\item 方法

\item 目的

\item 单桩竖向抗压静载试验

\item （

\item ）确定单桩竖向抗压极限承载力；

\item ）判定竖向抗压承载力是否满足设计要求；

\item ）通过桩身应变、位移测试，测定桩侧、桩端阻力，验证高应变法的单桩竖向抗压承载力检测结果。

\item 单桩竖向抗拔静载试验

\item ）确定单桩竖向抗拔极限承载力；

\item ）判断竖向抗拔承载力是否满足设计要求；

\item ）通过桩身应变、位移测试，测定桩的抗拔侧阻力。

\end{enumerate}

\newpage

\section*{24. 模板工程}

\addcontentsline{toc}{section}{24. 模板工程}

\begin{mainbox}

\textbf{提要：} 以下内容按课件自动提炼，适合快速背诵。

\end{mainbox}

\begin{enumerate}[label=\arabic*.]

\item 二、模板工程

\item （一）模板工程概述

\item 概述

\item （

\item ）模板——直接影响到混凝土成型的质量；

\item ）支架系统——直接影响到其他施工的安全。

\item （二）常见模板及其特性

\item 胶合板模板

\item 组合钢模板

\item 钢框木（竹）胶合板模板

\item 大模板

\item 组合铝合金模板

\item 滑升模板

\item 爬升模板

\item 飞模

\item 模壳模板

\item 胎模

\item 永久性压型钢板模板

\item 隧道模板

\item 早拆模板体系

\item 所用胶合板为高耐气候、耐水性的

\item I

\item 类木胶合板或竹胶合板。优点是自重轻、板幅大、板面平整、施工安装方便简单等。

\item 主要由钢模板、连接体和支撑体三部分组成，优点是轻便灵活、拆装方便、通用性强、周转率高等；缺点是接缝多且严密性差，导致混凝土成型后外观质量差。

\end{enumerate}

\newpage

\section*{25. 流砂及其防治}

\addcontentsline{toc}{section}{25. 流砂及其防治}

\begin{mainbox}

\textbf{提要：} 以下内容按课件自动提炼，适合快速背诵。

\end{mainbox}

\begin{enumerate}[label=\arabic*.]

\item 五、流砂及其防治

\item （一）流砂产生的条件

\item （二）流砂形成原因

\item 流砂形成原因

\item 内因

\item 土质问题——

\item 均匀：颗粒级配中，土的不均匀系数小于

\item ；

\item 粘性差 ：土的颗粒组成中，粘粒含量

\item <10

\item ％，粉粒含量＞

\item 75\%

\item 土松：天然空隙比

\item >0.75

\item 水多：含水量

\item >30\%

\item 因此，流砂易在粉土、细砂、粉砂和淤泥土中发生。

\item 外因

\item 与土重作用相反的动水压力——

\item 动水压力的性质 ：

\item （

\item ）作用方向与水流方向相同；

\item ）与水力梯度即水头差成正比；

\item ）与渗透路径成反比。

\end{enumerate}

\newpage

\section*{26. 混凝土小型空心砌块砌体工程}

\addcontentsline{toc}{section}{26. 混凝土小型空心砌块砌体工程}

\begin{mainbox}

\textbf{提要：} 以下内容按课件自动提炼，适合快速背诵。

\end{mainbox}

\begin{enumerate}[label=\arabic*.]

\item 三、混凝土小型空心砌块砌体工程

\item 混凝土小型空心砌块砌体工程

\item （

\item ）普通混凝土小型空心砌块砌体，不需对小砌块浇水湿润；如遇天气干燥炎热，宜在砌筑前对其喷水湿润；

\item ）轻骨料混凝土小砌块，应提前浇水湿润，块体的相对含水率宜为

\item 40\%

\item ～

\item 50\%

\item 。雨天及小砌块表面有浮水时，不得施工。

\item ）小砌块施工时，必须与砖砌体施工一样设立皮数杆、拉水准线。

\item ）小砌块墙体应孔对孔、肋对肋错缝搭砌。单排孔小砌块的搭接长度应为块体长度的

\item 1/2

\item ；多排孔小砌块的搭接长度可适当调整，但不宜＜小砌块长度的

\item 1/3

\item ，且不应＜

\item 90mm

\item 。墙体的个别部位不能满足上述要求时，应在灰缝中设置拉结钢筋或钢筋网片，但竖向通缝仍不得超过两皮小砌块。

\item ）砌筑应从转角或定位处开始，内外墙同时砌筑，纵横交错搭接。外墙转角处应使小砌块隔皮露端面；

\item T

\item 形交接处应使横墙小砌块隔皮露端面。

\item ）小砌块施工应对孔错缝搭砌，灰缝应横平竖直，宽度宜

\item 12mm

\item 。砌体水平、竖向灰缝的砂浆饱满度，按净面积计算不得低于

\item 90\%

\end{enumerate}

\newpage

\section*{27. 混凝土工程}

\addcontentsline{toc}{section}{27. 混凝土工程}

\begin{mainbox}

\textbf{提要：} 以下内容按课件自动提炼，适合快速背诵。

\end{mainbox}

\begin{enumerate}[label=\arabic*.]

\item 三、混凝土工程

\item （一）水泥的性能和应用

\item 概念

\item 无机水硬性胶凝材料；

\item 分类

\item 水硬性物质

\item 硅酸盐水泥、铝酸盐水泥、硫铝酸盐水泥、氟铝酸盐水泥、磷酸盐水泥等；

\item 用途性能

\item 通用水泥、特种水泥；

\item 混合材料品种、掺量

\item 硅酸盐水泥、普通硅酸盐水泥、矿渣硅酸盐水泥、火山灰质硅酸盐水泥、粉煤灰硅酸盐水泥、复合硅酸盐水泥；

\item 、常用水泥的技术要求

\item 凝结时间

\item 初凝

\item （

\item ）从水泥加水拌合起至水泥浆开始失去可塑性所需的时间；

\item ）六大常用水泥的初凝时间均不得短于

\item 45min

\item ；

\item 终凝

\item ）从水泥加水拌合起至水泥浆完全失去可塑性并开始产生强度所需的时间。

\item ）硅酸盐水泥的终凝时间不得长于

\item 6.5h

\item ）其他五类常用水泥的终凝时间不得长于

\end{enumerate}

\newpage

\section*{28. 砌筑工程季节性施工}

\addcontentsline{toc}{section}{28. 砌筑工程季节性施工}

\begin{mainbox}

\textbf{提要：} 以下内容按课件自动提炼，适合快速背诵。

\end{mainbox}

\begin{enumerate}[label=\arabic*.]

\item 五、砌体工程冬期施工

\item 砌体工程

\item （

\item ）冬期施工所用材料应符合下列规定：

\item ）砖、砌块在砌筑前，应清除表面污物、冰雪等，不得使用遭水浸和受冻后表面结冰、污染的砖或砌块；

\item ）砌筑砂浆宜采用普通硅酸盐水泥配制，不得使用无水泥拌制的砂浆；

\item ）现场拌制砂浆所用砂中不得含有直径大于

\item 10mm

\item 的冻结块或冰块；

\item ）石灰膏、电石渣膏等材料应有保温措施，遭冻结时应经融化后方可使用；

\item ）砂浆拌合水温不宜超过

\item ℃，砂加热温度不宜超过

\item ℃，且水泥不得与

\item ℃以上热水直接接触；砂浆稠度宜较常温适当增大，且不得二次加水调整砂浆和易性。

\item ）砌筑施工时，砂浆温度不应低于

\item ℃。当设计无要求，且最低气温等于或低

\item -15

\item ℃时，砌体砂浆强度等级应较常温施工提高一级。

\item ）砌体采用氯盐砂浆施工，每日砌筑高度不宜超过

\item 1.2m

\item ，墙体留置的洞口，距交接墙处不应小于

\item 500mm

\item 。

\item ）下列情况不得采用掺氯盐的砂浆砌筑砌体：

\end{enumerate}

\newpage

\section*{29. 砌筑砂浆}

\addcontentsline{toc}{section}{29. 砌筑砂浆}

\begin{mainbox}

\textbf{提要：} 以下内容按课件自动提炼，适合快速背诵。

\end{mainbox}

\begin{enumerate}[label=\arabic*.]

\item 一、砌筑砂浆

\item （一）原材料要求

\item 原材料要求

\item 水泥

\item （

\item ）进场使用前应有出厂合格证和复试合格报告；

\item ）强度等级——根据砂浆品种及强度等级要求选择；

\item ）$\leq$

\item M15

\item ——宜用

\item 32.5

\item 级的通用硅酸盐水泥或砌筑水泥；

\item ）＞

\item 42.5

\item 级普通硅酸盐水泥。

\item 砂

\item 宜用中砂，其中毛石砌体宜用粗砂。

\item 石灰膏

\item ）熟化时间——至少

\item 7d

\item （磨细——不少于

\item 2d

\item ）；

\item ）水泥石灰砂浆——不得采用脱水硬化的石灰膏。

\end{enumerate}

\newpage

\section*{30. 砖砌体工程}

\addcontentsline{toc}{section}{30. 砖砌体工程}

\begin{mainbox}

\textbf{提要：} 以下内容按课件自动提炼，适合快速背诵。

\end{mainbox}

\begin{enumerate}[label=\arabic*.]

\item 二、砖砌体工程

\item （一）砌块种类

\item 粘土多孔砖

\item 页岩空心砖

\item 煤矸石多孔砖

\item 页岩多孔砖

\item 普通砖

\item 混凝土空心砌块

\item 填有保温材料的砌块

\item 加气混凝土砌块

\item 陶粒混凝土空心砌块

\item 嵌锁型空心砌块

\item 砖

\item 烧结砖

\item 烧结普通砖

\item 240mm×115mm×53mm

\item ；

\item 烧结多孔砖

\item （

\item ）孔洞率不大于

\item 33\%

\item ，孔的尺寸小而数量多；

\item ）主要用于承重部位，砌筑时孔洞垂直于受压面。

\item 烧结空心砖

\end{enumerate}

\newpage

\section*{31. 网络优化}

\addcontentsline{toc}{section}{31. 网络优化}

\begin{mainbox}

\textbf{提要：} 以下内容按课件自动提炼，适合快速背诵。

\end{mainbox}

\begin{enumerate}[label=\arabic*.]

\item 二 、网络计划优化

\item 网络计划表示的逻辑关系

\item （

\item ）工艺关系：由工艺技术要求的工作先后顺序关系；

\item ）组织关系：施工组织时按需要进行的工作先后顺序安排。

\item ）通常情况下，网络计划优化时，只能调整工作间的组织关系。

\item 网络计划的优化目标

\item ）按计划任务的需要和条件可分为三方面：工期目标、费用目标和资源目标。

\item ）根据优化目标的不同，网络计划的优化相应分为工期优化、费用优化和资源优化三种。

\item （一）工期优化

\item 工期优化

\item ）工期优化：当网络计划的计算工期不满足要求工期时，通过压缩关键工作的持续时间以满足要求工期目标的过程。

\item ）网络计划工期优化的基本方法是在不改变网络计划中各项工作之间逻辑关系的前提下，通过压缩关键工作的持续时间来达到优化目标。

\item ）在工期优化过程中，按照经济合理的原则，不能将关键工作压缩成非关键工作。

\item ）此外，当工期优化过程中出现多条关键线路时，必须将各条关键线路的总持续时间压缩相同数值，否则，不能有效地缩短工期。

\item ）选择压缩对象时应考虑下列因素：

\item ）缩短持续时间对质量和安全影响不大的工作。

\item ）有充足备用资源的工作。

\item ）缩短持续时间所需增加的费用最少的工作。

\item ）选择关键工作压缩其持续时间时，应选择优选系数最小的关键工作。

\item 具体步骤

\item ）确定初始网络计划的计算工期和关键线路。

\item ）按要求工期计算应缩短的时间：$\triangle$

\item T

\end{enumerate}

\newpage

\section*{32. 脚手架工程的安全技术要求}

\addcontentsline{toc}{section}{32. 脚手架工程的安全技术要求}

\begin{mainbox}

\textbf{提要：} 以下内容按课件自动提炼，适合快速背诵。

\end{mainbox}

\begin{enumerate}[label=\arabic*.]

\item 三、脚手架工程的安全技术要求

\item （一）脚手架的施工准备工作

\item 施工准备工作

\item （

\item ）脚手架搭设之前，应根据工程的特点和施工工艺要求确定搭设与拆除的施工方案。

\item ）施工方案内容主要应包括：

\item ）材料要求。

\item ）基础要求。

\item ）荷载计算、计算简图、计算结果、安全系数。

\item ）立杆横距、立杆纵距、杆件连接、步距、允许搭设高度、连墙杆做法、门洞处理、剪刀撑要求、脚手板、挡脚板、扫地杆等构造要求。

\item ）脚手架搭设、拆除，安全技术措施及安全管理、维护、保养，以及平面图、剖面图、立面图、节点图要反映杆件连接、拉结基础等情况。

\item ）悬挑式脚手架有关悬挑梁、横梁等的加

\item 工

\item 节点图，悬挑梁与结构的连接节点，钢梁平面图，悬挑设计节点图。

\item （二）脚手架材质、构配件要求

\item 材料与规格

\item ）木制材料：木杆常用剥皮杉杆或落叶松，脚手板的厚度一般不小于

\item 50mm

\item ；

\item ）竹制材料：

\item 3\textasciitilde{}4

\item 年以上生长的毛竹，主要受力杆件的使用期限不宜超过

\item 年。

\item ）钢管：

\end{enumerate}

\newpage

\section*{33. 脚手架的种类}

\addcontentsline{toc}{section}{33. 脚手架的种类}

\begin{mainbox}

\textbf{提要：} 以下内容按课件自动提炼，适合快速背诵。

\end{mainbox}

\begin{enumerate}[label=\arabic*.]

\item 一、脚手架的种类

\item 概念

\item 它是为保证高处作业安全、顺利进行施工而搭设的工作平台或作业通道。在结构施工、装修施工和设备管道的安装施工中，都需要按照操作要求搭设脚手架。

\item 分类

\item 按其搭设位置分

\item 外脚手架和里脚手架；

\item 按其所用材料分

\item 木脚手架、竹脚手架与金属脚手架；

\item 按其构造形式分

\item 多立杆式、框式、桥式、吊式、挂式、升降式以及用于层间操作的工具式脚手架；

\item 按搭设高度分

\item 高层脚手架和普通脚手架；

\item 基本要求

\item 其宽度应满足工人操作、材料堆置和运输的需要；坚固稳定；装拆简便；能多次周转使用。

\item 扣件式钢管脚手架

\item 用扣件钢管搭设的悬挑外脚手架

\item 扣件式钢管脚手架组成

\item 扣件式钢管脚手架立面图

\item 扣件式钢管脚手架剖面图

\item 旋转扣件

\item 直角扣件

\item 对接

\item 扣件

\item 冲压钢脚手板

\end{enumerate}

\newpage

\section*{34. 边坡开挖}

\addcontentsline{toc}{section}{34. 边坡开挖}

\begin{mainbox}

\textbf{提要：} 以下内容按课件自动提炼，适合快速背诵。

\end{mainbox}

\begin{enumerate}[label=\arabic*.]

\item 一、边坡开挖

\item 土方边坡坡度

\item =H/B=1/(B/H)=1/m

\item 坡度系数

\item m=B/H

\item 边坡稳定的条件：

\item T——

\item 土体下滑力。下滑土体的分力，受坡上荷载、雨水、静水压力影响。

\item C——

\item 土体抗剪力。由土质决定，受气候、含水量及动水压力影响。

\item 《

\item 建筑地基基础工程施工质量验收标准

\item 》

\item （

\item GB 50202-2018

\item ），在边坡整体稳定的情况下，

\item 临时性挖方工程的边坡坡率允许值

\item 应符合下表的规定或经设计计算确定。

\item 放坡开挖施工应符合下列要求

\item ）应按先降低地下水位，然后开挖，再做坡面护理的工序进行施工；

\item ）开挖前应校核开挖尺寸线，检查地面排水措施和降水场地的水位标高，符合要求后方可开挖；

\item ）土方开挖应按先上后下的开挖顺序，分段、分层按设计要求开挖，分层、分段开挖尺寸应符合设计工况要求，开挖过程中应确保坡壁无超挖，坡面无虚土，坡面坡度与平整度应符合设计要求；

\item ）黏性土分段开挖长度宜取

\item 10m\textasciitilde{}15m

\end{enumerate}

\newpage

\section*{35. 钢筋工程}

\addcontentsline{toc}{section}{35. 钢筋工程}

\begin{mainbox}

\textbf{提要：} 以下内容按课件自动提炼，适合快速背诵。

\end{mainbox}

\begin{enumerate}[label=\arabic*.]

\item 一、钢筋工程

\item （一）建筑钢材的主要钢种

\item 建筑钢材：钢结构用钢、钢筋混凝土结构用钢、建筑装饰用钢制品。

\item 主要钢种

\item 化学成分

\item 碳素钢——低碳、

\item 中碳

\item \relax{}[0.25\%

\item ～

\item 0.6\%]

\item 、高碳；

\item 合金钢——低合金、

\item 中合金

\item \relax{}[5\%

\item 10\%]

\item 、高合金；

\item 有害杂质

\item 普通钢、优质钢、高级优质钢、特级优质钢；

\item 碳素结构钢

\item （

\item ）

\item Q235

\item －

\item AF

\end{enumerate}

\newpage

\section*{36. 钢筋混凝土灌注桩}

\addcontentsline{toc}{section}{36. 钢筋混凝土灌注桩}

\begin{mainbox}

\textbf{提要：} 以下内容按课件自动提炼，适合快速背诵。

\end{mainbox}

\begin{enumerate}[label=\arabic*.]

\item 四、钢筋混凝土灌注桩

\item （一）泥浆护壁灌注桩

\item 成孔工艺

\item 正（反）循环钻机、冲击钻机、旋挖钻机、多支盘灌注桩机、扩底机械钻具等

\item 工艺流程

\item 场地平整→桩位放线→开挖浆池、浆沟→护筒埋设→钻机就位、孔位校正→成孔、泥浆循环、清除废浆、泥渣→清孔换浆→终孔验收→下钢筋笼和钢导管→二次清孔

\item →

\item 清孔质量检查

\item 浇筑水下混凝土→成桩。

\item 施工要求

\item （

\item ）试孔——工艺性试成孔，数量$\geq$

\item 根。

\item ）护壁泥浆——可采用原土造浆，不适用的土层应制备泥浆。施工时，钻孔内泥浆液面高出地下水位

\item 0.5m

\item 。

\item ）成孔——

\item 正、反循环成孔机具——桩型、地质条件及成孔工艺选择，砂土层成孔宜选用反循环钻机；

\item 冲击钻成孔——遇岩石表面不平或遇孤石时，应向孔内投入黏土、块石，将孔底表面填平后低锤快击，形成挤密平台，再进行正常冲击。

\item 多支盘灌注桩成孔——可采用泥浆护壁成孔、干作业成孔、水泥注浆护壁成孔、重锤捣扩成孔等方法。

\item ）清孔——正循环、泵吸反循环、气举反循环等。

\item 沉渣厚度——端承桩$\leq$

\item 50mm

\item ，摩擦桩$\leq$

\end{enumerate}

\newpage

\section*{37. 钢筋混凝土预制桩}

\addcontentsline{toc}{section}{37. 钢筋混凝土预制桩}

\begin{mainbox}

\textbf{提要：} 以下内容按课件自动提炼，适合快速背诵。

\end{mainbox}

\begin{enumerate}[label=\arabic*.]

\item 三、钢筋混凝土预制桩

\item （一）锤击沉桩法

\item 施工程序

\item 确定桩位和沉桩顺序→桩机就位→吊桩喂桩→校正→锤击沉桩→接桩→再锤击沉桩→送桩→收锤→切割桩头。

\item 施工要求

\item （

\item ）混凝土强度达到

\item 70\%

\item 后方可起吊，达到

\item 100\%

\item 后方可运输和打桩。

\item ）单节桩采用两支点起吊时，吊点距桩端宜为

\item 0.2L

\item （桩段长）。吊运过程中严禁采用拖拉取桩方法。

\item ）接桩——接头宜高出地面

\item 0.5

\item ～

\item 1m

\item 。

\item 接桩方法——焊接、螺纹接头和机械啮合接头等。

\item ）沉桩顺序——先深后浅、先大后小、先长后短、先密后疏的次序进行。对于密集桩群应控制沉桩速率，宜从中间向四周或两边对称施打；当一侧毗邻建筑物时，由毗邻建筑物处向另一方向施打。

\item ）终沉标准

\item 终止沉桩——桩端标高控制为主，贯入度控制为辅；当桩终端达到坚硬，硬塑黏性土，中密以上粉土、砂土、碎石土及风化岩时，可以贯入度控制为主，桩端标高控制为辅；贯入度达到设计要求而桩端标高未达到时，应继续锤击

\item 阵，按每阵

\end{enumerate}

\newpage

\section*{38. 非（无）节奏流水施工}

\addcontentsline{toc}{section}{38. 非（无）节奏流水施工}

\begin{mainbox}

\textbf{提要：} 以下内容按课件自动提炼，适合快速背诵。

\end{mainbox}

\begin{enumerate}[label=\arabic*.]

\item 、无节奏流水施工

\item （

\item ）概念：

\item 同一施工过程在各施工段上的流水节拍不完全相等

\item ，

\item 不同的施工过程之间流水节拍也不完全相等

\item ，在这样的条件下组织施工的方式称为无节奏流水施工。

\item ）基本特征：

\item ）主要参数的确定：

\item 例题：

\item 某工程有

\item A

\item 、

\item B

\item C

\item D

\item 四个施工过程，平面上划分成四个施工段，每个施工过程在各个施工段上的流水节拍如下表所示，试计算工期并绘制流水施工进度计划。

\item 解：

\item ）已知：施工过程

\item n=4

\item ，施工段

\item m=4

\item ；

\item ）确定流水步距：“累加数列法”

\end{enumerate}

\newpage

\section*{39. 预应力施加方法}

\addcontentsline{toc}{section}{39. 预应力施加方法}

\begin{mainbox}

\textbf{提要：} 以下内容按课件自动提炼，适合快速背诵。

\end{mainbox}

\begin{enumerate}[label=\arabic*.]

\item 六、预应力施加方法

\item 分类

\item 预应力的施加方法，根据与构件制作相比较的先后顺序，分为先张法、后张法两大类。

\item 特点

\item （

\item ）先张法——先张拉预应力筋，再浇筑砼；预应力是靠预应力筋与砼之间的粘结力传递给砼，并使其产生预压应力。

\item ）后张法——先浇筑砼，达到一定强度后，再在其上张拉预应力筋；预应力是靠锚具传递给砼，并使其产生预压应力。

\item 锚具作为预应力构件的一个组成部分，永远留在构件上，不能重复使用。

\item （一）先张法施工

\item 先张法工艺流程

\item ——

\item 预应力筋的张拉

\item ）施加预应力宜采用一端张拉工艺，张拉控制应力和程序按设计要求进行。

\item ）当设计无具体要求时，一般采用：

\item →

\item 1.03

\item σ

\item con

\item 或

\item 1.05

\item （持荷

\item 2min

\item ）→σ

\item ，目的是为了减少预应力的松弛损失。

\end{enumerate}

\newpage

\section*{40. 预应力混凝土概述}

\addcontentsline{toc}{section}{40. 预应力混凝土概述}

\begin{mainbox}

\textbf{提要：} 以下内容按课件自动提炼，适合快速背诵。

\end{mainbox}

\begin{enumerate}[label=\arabic*.]

\item 一、预应力混凝土概述

\item （一）预应力混凝土工作原理

\item （二）预应力混凝土的特点

\item 优点

\item 抗裂性好，刚度大；

\item 节省材料，减小自重；

\item 可以减小混凝土梁的竖向剪力和主拉应力；

\item 提高受压构件的稳定性；

\item 提高构件的耐疲劳性能；

\item 预应力可以作为结构构件连接的手段，促进大跨结构新体系与施工方法的发展。

\item 缺点

\item 工艺较复杂，对质量要求高；

\item 需要有一定的专门设备，如张拉机具、灌浆设备等；

\item 预应力混凝土结构的开工费用较大，对构件数量少的工程成本较高；

\item 预应力反拱度不易控制。

\item 二、预应力钢筋的种类

\item 冷拉钢筋

\item 冷拉钢筋的塑性和弹性模量有所降低而屈服强度和硬度有所提高，可直接用作预应力钢筋。

\item 高强度钢丝

\item 常用的高强钢丝分为冷拉和矫直回火两种，按外形分为光面、刻痕和螺旋肋三种。还有消除应力钢丝，低松弛钢丝。

\item 钢绞线

\item 股钢绞线由于面积较大、柔软、施工定位方便，适用于先张法和后张法预应力结构与构件，是目前国内外应用最广的一种预应力筋。

\item 热处理钢筋

\item 具有高强度、高韧性和高黏结力等优点，直径为

\end{enumerate}

\newpage

\end{document}